\documentclass[10pt]{article}
\usepackage[margin=0.76in]{geometry}
\usepackage{amsmath,amssymb,amsthm,mathtools}
\usepackage{graphicx}
\usepackage{xcolor}
\usepackage{microtype}
\usepackage[hidelinks]{hyperref}
\setlength{\parindent}{0pt}
\setlength{\parskip}{5pt}
\newtheorem{theorem}{Theorem}
\newtheorem{lemma}[theorem]{Lemma}
\newtheorem{corollary}[theorem]{Corollary}
\newcommand{\R}{\mathbb R}
\newcommand{\Mor}{\mathcal L}
\newcommand{\WMor}{W\mathcal L}
\newcommand{\avgint}{\mathop{\rlap{\raisebox{0.18ex}{--}}\!\int}}

\title{Distance-to-Subspace Weights Disprove the Proposed\\
Muckenhoupt--Reverse-Hölder Sufficiency\\
\large A full negative answer to the final question of arXiv:1910.13285}
\author{Automated proof attempt\\
\texttt{agent\_lane\_03} \quad model: GPT5.6}
\date{11 August 2026}

\begin{document}
\maketitle

\begin{center}
\fcolorbox{red!55!black}{red!4}{\parbox{0.92\linewidth}{
\textbf{Status: full counterexample, likely valid, pending human verification.}
The proposed sufficiency of
$A_{p+\lambda/n}\cap RH_{n/(n-\lambda)}$ is false.  In fact, a family of
distance-to-subspace power weights in that class makes the Hardy--Littlewood
maximal operator fail even from strong to weak weighted Morrey space.}}
\end{center}

\section*{Source question}

For $1\le p<\infty$ and $0<\lambda<n$, the source defines
\[
 \|f\|_{\Mor^{p,\lambda}(w)}
 =\sup_{x\in\R^n,r>0}
 \left(r^{-\lambda}\int_{B(x,r)}|f|^p w\right)^{1/p},
\]
and
\[
 \|g\|_{\WMor^{p,\lambda}(w)}
 =\sup_{x,r,a>0}
 \left(r^{-\lambda}a^p
 w\{u\in B(x,r):|g(u)|>a\}\right)^{1/p}.
\]
After establishing shifted sharp ranges for point powers, the authors ask
whether the whole class
\begin{equation}\label{eq:asked-class}
 A_{p+\lambda/n}\cap RH_{n/(n-\lambda)}
\end{equation}
is sufficient in their Theorem 1.1 \cite{source}.  Since the pair
$(|f|,Mf)$ satisfies the theorem's weighted hypothesis, this would in
particular imply Morrey boundedness of the Hardy--Littlewood maximal operator
for every weight in \eqref{eq:asked-class}.  The appended crop reproduces the
question exactly from page 16.

\section*{Proof intuition}

The point powers treated in the source vanish at a zero-dimensional set.  A
ball larger than that small low-weight region pays for spreading its mass in
every direction.  The proposed weight classes do not encode this geometry.

Replace the point by a positive-dimensional subspace and use a power of the
transverse distance.  A thin tube is very cheap in weighted Morrey norm
because the weight vanishes across its transverse thickness.  Its ordinary
Lebesgue mass is only one power of that thickness smaller, and the maximal
operator spreads this mass over a fixed ball.  If the subspace has enough
tangential dimension to absorb the Morrey normalization, the two scalings
are incompatible.  The strict gap in the $A_{p+\lambda/n}$ range supplies
exactly the exponents needed for the contradiction.

\section*{A weight-class calculation}

We use the following standard subspace-power fact, including a short proof to
fix all endpoint conventions.

\begin{lemma}[Powers of transverse distance]\label{lem:weight}
Write $\R^n=\R^d\times\R^m$, where $d+m=n$ and $m\ge1$, and let
\(w(y,z)=|z|^\beta\).  For $q>1$,
\[
 w\in A_q(\R^n)\quad\Longleftrightarrow\quad
 -m<\beta<m(q-1).
\]
If $\beta\ge0$, then $w\in RH_s(\R^n)$ for every finite $s>1$.
\end{lemma}

\begin{proof}
For every $\gamma>-m$ and every ball
$B=B((y_0,z_0),r)\subset\R^d\times\R^m$, scaling the transverse variable
in the section integral gives
\begin{equation}\label{eq:power-average}
 \frac1{|B|}\int_B |z|^\gamma\,dy\,dz
 \asymp (|z_0|+r)^\gamma.
\end{equation}
Indeed, when $|z_0|\ge2r$ the integrand is pointwise comparable to
$|z_0|^\gamma$.  When $|z_0|<2r$, scale by $r$; the remaining integrals have
translation parameter in a compact set and are uniformly positive and
finite precisely for $\gamma>-m$.

Apply \eqref{eq:power-average} first with $\gamma=\beta$ and then with
$\gamma=-\beta/(q-1)$.  Their powers cancel in the $A_q$ product, and the
second integral is locally finite exactly when $\beta<m(q-1)$; local
integrability of $w$ itself requires $\beta>-m$.  For $RH_s$, apply
\eqref{eq:power-average} with $\gamma=\beta s$ and $\gamma=\beta$.
When $\beta\ge0$ both are admissible and the two sides of the reverse
Hölder inequality are comparable to $(|z_0|+r)^\beta$.
\end{proof}

\section*{The counterexample family}

\begin{theorem}[Failure throughout a strict parameter strip]\label{thm:main}
Let $1\le p<\infty$, $0<\lambda<n$, and write
$\R^n=\R^d\times\R^m$ with
\[
 d+m=n,\qquad m\ge1,\qquad d\ge\lambda.
\]
Choose
\begin{equation}\label{eq:beta-range}
 m(p-1)<\beta<m\left(p-1+\frac{\lambda}{n}\right),
 \qquad w(y,z)=|z|^\beta.
\end{equation}
Then
\[
 w\in A_{p+\lambda/n}\cap RH_{n/(n-\lambda)},
\]
but the Hardy--Littlewood maximal operator is not bounded
\[
 M:\Mor^{p,\lambda}(w)\longrightarrow\WMor^{p,\lambda}(w).
\]
Consequently it is not bounded on $\Mor^{p,\lambda}(w)$ either.
\end{theorem}

\begin{proof}
Put $q=p+\lambda/n$.  The upper inequality in \eqref{eq:beta-range} is
$\beta<m(q-1)$, while the lower one gives $\beta>0$.  Lemma
\ref{lem:weight} therefore proves membership in both required classes.

For $0<t<1$, set
\[
 E_t=B_d(0,1)\times B_m(0,t),\qquad f_t=\chi_{E_t}.
\]
We first estimate the strong input norm.  Let $B=B((a,b),r)$ be any ball.
If $r\le t$, then on $B\cap E_t$ one has $w\le t^\beta$, and hence
\[
 r^{-\lambda}\int_B f_t^p w
 \le C r^{n-\lambda}t^\beta
 \le C t^{n-\lambda+\beta}
 \le C t^{m+\beta};
\]
the last inequality uses $n-\lambda=d+m-\lambda\ge m$ and $t<1$.
If $t<r\le1$, transverse integration and the bound by a tangential
$d$-ball of radius $r$ give
\[
 r^{-\lambda}\int_B f_t^p w
 \le C r^{d-\lambda}\int_{|z|<t}|z|^\beta dz
 \le C t^{m+\beta},
\]
because $d\ge\lambda$.  Finally, if $r\ge1$, use the total weighted mass
of $E_t$ and $r^{-\lambda}\le1$.  Uniformly in all cases,
\begin{equation}\label{eq:input}
 \|f_t\|_{\Mor^{p,\lambda}(w)}le C t^{(m+\beta)/p}.
\end{equation}

For every $x\in B_n(0,1)$, the ball $B_n(x,3)$ contains $E_t$.  Since
$|E_t|=c_{d,m}t^m$, the centered (and therefore also the usual)
Hardy--Littlewood maximal operator satisfies
\begin{equation}\label{eq:point-lower}
 Mf_t(x)\ge c t^m\qquad (x\in B_n(0,1)).
\end{equation}
Choose the weak-norm level $a=ct^m/2$ and the unit ball in its definition.
Because $0<w(B_n(0,1))<\infty$, \eqref{eq:point-lower} yields
\begin{equation}\label{eq:output}
 \|Mf_t\|_{\WMor^{p,\lambda}(w)}\ge c' t^m.
\end{equation}
Combining \eqref{eq:input} and \eqref{eq:output}, the operator-norm quotient
is at least
\[
 c''t^{,m-(m+\beta)/p}
 =c''t^{[m(p-1)-\beta]/p}\longrightarrow\infty
 \qquad(t\downarrow0),
\]
by the strict lower inequality in \eqref{eq:beta-range}.  This proves failure
even into weak Morrey space.
\end{proof}

\begin{corollary}[A concrete strong-type disproof]\label{cor:explicit}
In $\R^2$, take
\[
 p=2,\qquad\lambda=1,\qquad
 w(x_1,x_2)=|x_2|^{5/4}.
\]
Then $w\in A_{5/2}\cap RH_2$, but $M$ is not bounded on
$\Mor^{2,1}(w)$; indeed it is not bounded from that space to
$\WMor^{2,1}(w)$.
\end{corollary}

\begin{proof}
Use Theorem \ref{thm:main} with $d=m=1$ and $\beta=5/4$.  The two strict
conditions are
\[
 1=m(p-1)<\frac54<\frac32
 =m\left(p-1+\frac\lambda n\right).
\]
More explicitly, horizontal strips satisfy
\(|f_t|_{\Mor^{2,1}(w)}=O(t^{9/8})\), whereas
\(|Mf_t|_{\WMor^{2,1}(w)}\ge ct\); their quotient grows as
$t^{-1/8}$.
\end{proof}

\section*{Scope, verification, and literature boundary}

Corollary \ref{cor:explicit} alone fully answers the proposed universal
sufficiency negatively, without relying on the weak endpoint $p=1$.  Theorem
\ref{thm:main} shows that the obstruction is robust and even defeats the
strong-to-weak estimate.  The construction requires an integer codimension
$m\le n-\lambda$; it makes no claim about parameter regimes without such an
$m$, nor about a replacement characterization of all admissible weights.

The script \path{code/check_scaling.py} checks the exact rational inequalities
and the $t^{-1/8}$ quotient in the explicit example.  It is not used in the
proof.  The ball estimates above were also checked at the three scale regimes
$r\le t$, $t<r\le1$, and $r\ge1$; both boundary exponents match.

The bounded novelty audit searched the run indexes; exact class and source
phrases; title/author citations; and hyperplane-power, subspace-distance, and
weighted-Morrey maximal-operator terms through 11 August 2026.  No matching
resolution was found.  The authors' later intrinsic-weight paper
\cite{followup2020} gives necessary conditions and sufficiency with an extra
local $A_p$ hypothesis, while Lerner's 2022 note \cite{lerner2022} still
describes the broader global intrinsic characterization as open.  Neither
contains this construction.  This is a bounded audit, not an exhaustive
priority claim.

\begin{thebibliography}{9}
\bibitem{source}
J. Duoandikoetxea and M. Rosenthal,
\emph{Boundedness of operators on certain power-weighted Morrey spaces beyond
the Muckenhoupt weights}, Potential Anal. \textbf{53} (2020), 1255--1268;
arXiv:1910.13285.  Included as \texttt{source\_paper.pdf}.

\bibitem{followup2020}
J. Duoandikoetxea and M. Rosenthal,
\emph{Muckenhoupt-type conditions on weighted Morrey spaces},
J. Fourier Anal. Appl. \textbf{27} (2021), Article 32;
arXiv:2010.00250.  Included as \texttt{supporting\_2010.00250.pdf}.

\bibitem{lerner2022}
A. K. Lerner,
\emph{A note on the maximal operator on weighted Morrey spaces},
arXiv:2211.07974 (2022).  Included as
\texttt{supporting\_2211.07974.pdf}.
\end{thebibliography}

\clearpage
\section*{Exact source excerpt: final open question (page 16)}
\begin{center}
\includegraphics[width=0.96\linewidth,keepaspectratio]{figures/open_question_crop.png}
\end{center}

\end{document}
